\documentclass[11pt]{article}
\usepackage{amsmath,amssymb,amsthm}
\newtheorem{proposition}{Proposition}
\title{A denominator obstruction for five-row local models}
\author{Research note; independently reviewed}
\date{30 September 2026}
\begin{document}
\maketitle
\begin{proposition}
A rational five-row, four-column comparison model in centered
coordinates $x=2p-1$ cannot have all its entries in $\mathbb Z_2$. Thus at least one entry has even denominator
in lowest terms. No real monotonicity assumption is needed.
\end{proposition}
\begin{proof}
Suppose the entries $x_{j,i}$ are 2-adic integers. The nonempty
square-free coordinate sums must be zero in odd degree, $5/3$ in
degree two, and $1$ in degree four. Modulo two these moments, together
with the empty sum $5$, determine the parity of the multiplicity of
each row pattern in $\mathbb F_2^4$. Boolean inclusion--exclusion gives
odd multiplicity for exactly the five patterns
\[
 (1,1,1,1),\ (0,1,1,1),\ (1,0,1,1),\ (1,1,0,1),\ (1,1,1,0).
\]
For completeness, the parity for a pattern whose support has size $k$
is $\sum_{s=k}^4\binom{4-k}{s-k}m_s$, where
$(m_0,m_1,m_2,m_3,m_4)=(1,0,1,0,1)$; it equals one precisely for
$k=3,4$. Since there are only five rows, each displayed pattern occurs
exactly once.

Label the all-ones row by $*$ and the row missing coordinate $i$ by
$i$. Write $x_{j,t}=b_{j,t}+2y_{j,t}$ with the displayed Boolean base
matrix $b$, and put $a_{ji}=y_{j,i}\pmod2$. The first moments modulo
four give $\sum_t y_{j,t}=0\pmod2$, since $\sum_t b_{j,t}=4$.
For $j\ne k$, the pair base sum is three, and the required pair sum
$5/3$ also equals three modulo four. Its first-order expansion
therefore gives
\[
 a_{jk}+a_{kj}=0\pmod2.
\]
For a triple of distinct columns $j,k,l$, the Boolean base sum is two.
The first-order correction, divided by two and taken modulo two, is
\[
 (a_{jk}+a_{kj})+(a_{jl}+a_{lj})+(a_{kl}+a_{lk})=0.
\]
The triple sum is consequently two modulo four, whereas it is
required to be zero. This contradiction proves the assertion.
\end{proof}
The proposition is only a denominator restriction. It does not
exclude points over $\mathbb Q_2$ with negative valuations.
\end{document}
